% !TeX root = ../../Book2.tex
\section{Schwartz 空间}
\label{b2:sec:0501}

紧支集光滑函数适合局部微分，但 Fourier 变换通常不保持紧支集。另一方面，反演与分部积分希望函数及其全部导数都足够快地下降。Schwartz 空间同时满足这些要求，使变换、微分和多项式乘法可以在同一个测试空间内反复使用。

\subsection{快速下降的光滑函数}
\label{b2:subsec:050101}

\begin{definition}\label{b2:def:05-schwartz}
设 \(\varphi\in C^\infty(\mathbb R^d)\)。若对每对多重指标 \(\alpha,\beta\) 都有
\[
 p_{\alpha,\beta}(\varphi)\coloneqq
 \sup_{x\in\mathbb R^d}|x^\alpha\partial^\beta\varphi(x)|<\infty,
\]
则称 \(\varphi\) 为\term[schwartz-hanshu]{Schwartz 函数}[Schwartz function]。全部这样的函数组成\term[schwartz-kongjian]{Schwartz 空间}[Schwartz space]，记为 \(\mathcal S(\mathbb R^d)\)。
\end{definition}
\symidx{\(\mathcal S(\mathbb R^d)\)：Schwartz 测试函数空间}

定义控制的是每一阶导数，不只是函数本身。Gauss 函数及光滑紧支集函数都属于 \(\mathcal S\)。函数 \((1+|x|^2)^{-M}\) 虽光滑并有较快衰减，却对足够大的多项式权重失去有界性，所以对任意固定 \(M\) 都不属于 \(\mathcal S\)。

\subsection{Schwartz 半范数}
\label{b2:subsec:050102}

为统一记号，也可使用递增半范数
\[
 p_m(\varphi)\coloneqq
 \max_{|\beta|\leq m}\sup_x(1+|x|)^m|\partial^\beta\varphi(x)|,
 \quad m=0,1,2,\ldots.
\]
它们与前述半范数族给出相同的收敛：单项式由径向权控制；反过来，\((1+|x|)^m\) 由有限多个不超过 \(m\) 阶的绝对单项式之和控制。规定 \(\varphi_n\to\varphi\) 于 \(\mathcal S\)，意为每个 \(m\) 都有 \(p_m(\varphi_n-\varphi)\to0\)。
\symidx{\(p_m\)：Schwartz 空间的递增半范数}

\begin{proposition}\label{b2:prop:05-schwartz-complete}
上述收敛由平移不变度量
\[
 \mathrm{d}_{\mathcal S}(\varphi,\psi)
 =\sum_{m=0}^\infty2^{-m-1}\min\{1,p_m(\varphi-\psi)\}
\]
刻画，且 \(\mathcal S\) 对此度量完备。对每个有限 \(p\)，包含映射 \(\mathcal S\subset L^p\) 连续；对 \(p=\infty\) 也连续。
\end{proposition}
\begin{proof}
每个半范数趋零时，先截去度量级数的统一小尾部，再控制有限头部，便得度量趋零。反过来，度量趋零使每个固定非负项趋零，因而对应半范数趋零。同样论证说明度量 Cauchy 等价于逐半范数 Cauchy。

设 \((\varphi_n)\) 为这样的 Cauchy 列。每个导数 \(\partial^\beta\varphi_n\) 一致收敛到连续函数 \(v_\beta\)。在任意坐标线段上使用微积分基本定理并取一致极限，得到 \(\partial_jv_\beta=v_{\beta+e_j}\)。因此 \(v_0\) 光滑，全部导数为 \(v_\beta\)。对固定权重与导数阶，把 Cauchy 不等式逐点取极限再取上确界，得到 \(p_m(\varphi_n-v_0)\to0\)，且 \(p_m(v_0)<\infty\)。这证明完备性。

若选整数 \(m\) 使 \(mp>d\)，则 \(\|\varphi\|_p\leq p_m(\varphi)\|(1+|x|)^{-m}\|_p\)；无穷端点由 \(p_0\) 控制。故包含连续。
\end{proof}

完备可度量的局部凸空间称为\term[frechet-kongjian]{Fréchet 空间}[Fréchet space]，所以上述空间属于这一类。与 Banach 空间不同，此处的拓扑由一列兼顾越来越多导数和权重的半范数指定；本章只使用这些具体半范数，不依赖抽象局部凸空间理论。

\begin{proposition}\label{b2:prop:05-compact-cutoff}
空间 \(C_c^\infty(\mathbb R^d)\) 在 \(\mathcal S\) 中稠密。
\end{proposition}
\begin{proof}
取光滑截断 \(\chi\)，在单位球上为 \(1\)，在半径 \(2\) 的球外为零，令 \(\chi_R(x)=\chi(x/R)\)。对 \(\varphi\in\mathcal S\)，Leibniz 公式将 \(\partial^\beta((1-\chi_R)\varphi)\) 写成有限和。未落在截断上的导数项支撑于 \(|x|\geq R\)，乘 \((1+|x|)^m\) 后由 \(p_{m+1}(\varphi)/(1+R)\) 控制；落在截断上的每个导数另产生 \(R^{-1}\) 或更高负幂，也只支撑于这个外部区域。因此 \(p_m((1-\chi_R)\varphi)\to0\) 对每个 \(m\) 成立，而 \(\chi_R\varphi\in C_c^\infty\)。
\end{proof}

这里的稠密性与第一卷测试函数空间 \(\mathcal D\) 的收敛规则不同：截断半径随 \(R\) 增大，并不存在共同紧支集；所以该逼近是在 \(\mathcal S\) 中成立，不能称为 \(\mathcal D\) 中的收敛。

\subsection{微分}
\label{b2:subsec:050103}

对每个多重指标 \(\gamma\)，有 \(p_m(\partial^\gamma\varphi)\leq p_{m+|\gamma|}(\varphi)\)，故微分是 \(\mathcal S\) 上的连续线性映射。分部积分也没有无穷远边界项：可以先乘以 \(\chi_R\)，所有边界误差由快速衰减趋零，再令 \(R\to\infty\)。这个截断办法比单独写“边界项为零”更适合多维积分。

\subsection{多项式乘法}
\label{b2:subsec:050104}

\begin{proposition}\label{b2:prop:05-multiplication}
若 \(a\in C^\infty\)，且每个导数 \(\partial^\gamma a\) 都至多按某个多项式增长，则 \(\varphi\mapsto a\varphi\) 连续地映 \(\mathcal S\) 到自身。特别地，多项式乘法连续。
\end{proposition}
\begin{proof}
固定 \(m\)。Leibniz 公式只用到有限个 \(|\gamma|\leq m\) 的导数，所以存在共同的 \(N,C\)，使这些导数均不超过 \(C(1+|x|)^N\)。于是 \(p_m(a\varphi)\leq C_m\cdot p_{m+N}(\varphi)\)，即得连续性。\(N\) 可随 \(m\) 改变，不要求全部导数有同一个增长阶。
\end{proof}

如果只知道 \(a\) 本身有界而不控制其导数，这个结论便不能照用。例如高度振荡的光滑函数可能具有极大的导数；乘积的高阶半范数必须逐项检查。

\subsection{卷积}
\label{b2:subsec:050105}

\begin{proposition}\label{b2:prop:05-schwartz-convolution}
若 \(\varphi,\psi\in\mathcal S\)，则 \(\varphi*\psi\in\mathcal S\)，且卷积是连续双线性映射。并有 \(\partial^\beta(\varphi*\psi)=(\partial^\beta\varphi)*\psi\)。
\end{proposition}
\begin{proof}
导数公式由各阶导数的有界性与 \(\psi\) 的可积性逐次使用控制收敛。再利用 \(1+|x|\leq(1+|x-y|)(1+|y|)\)，对 \(|\beta|\leq m\) 得
\[
 (1+|x|)^m|\partial^\beta(\varphi*\psi)(x)|
 \leq p_m(\varphi)\int(1+|y|)^m|\psi(y)|\,\mathrm{d}y
 \leq C_d\cdot p_m(\varphi)\cdot p_{m+d+1}(\psi).
\]
对 \(x,\beta\) 取上确界得到半范数估计，也同时证明双线性连续性。
\end{proof}
